%% The abstract, input by paper.tex. Numbers are from loc.py (ISWIMμ, the
%% default report); keep them in sync.

Denotational semantics promises compositional, equational reasoning about
programs, yet mechanized compiler verification has been almost entirely
operational. The mechanized denotational proofs that exist stop at a
terminating source language, a first-order language, or a single pass,
because a Turing-complete, higher-order language requires a semantics that
accounts for divergence. We show that a denotational semantics based on
graph models, in which a denotation is a set of finite input/output tables,
meets this requirement. We prove correct, in Agda, a ten-pass compiler from
a call-by-value $\lambda$-calculus with iso-recursive types, which has
general recursion and divergence, through closure conversion, to a typed
assembly language. Every intermediate language has a graph-model
semantics, and most passes are proved by an equation between the
denotations of a program before and after the pass. The exception is the
step of closure conversion that lowers closures into code and environment
pairs: its reflection direction needs a new kind of logical relation
between two denotational semantics, the approximant logical relation,
defined by recursion on the depth of a finite approximation of the
target's meaning, because recursive types rule out recursion on types.
Adequacy at both ends yields an end-to-end theorem about running code: the
source and the compiled program return the same integer, or both diverge.
To measure the cost, we prove the same compiler correct in three
operational styles. The graph-model proofs of the passes total 1958 lines,
between those for clocked interpreters (1653) and abstract machines (2258)
and well below reduction semantics (3207); adequacy adds 1532. The graph
model gives the smallest proofs of inlining, hoisting and conversion to a
control-flow graph, and its cost concentrates in closure lowering.
